% Copyright 2026 Sewoong Lee
% 
% This work may be distributed and/or modified under the conditions of the
% LaTeX Project Public License, either version 1.3c of this license or (at your
% option) any later version. The latest version of this license is at:
% https://www.latex-project.org/lppl.txt
% 
% This work has the LPPL maintenance status `maintained'.
% The Current Maintainer is Sewoong Lee <samuel27@illinois.edu>.
% 
% The work consists of README.md, LICENSE, beamercolorthemechalkboard.sty,
% beamercolorthemechalkboard-doc.tex, and chalkboard-example.tex.
% The derived files are beamercolorthemechalkboard-doc.pdf and chalkboard-example.pdf.

\documentclass[11pt]{article}
\usepackage[T1]{fontenc}
\usepackage[margin=27mm]{geometry}
\usepackage{xcolor,booktabs}
\usepackage[colorlinks=true,urlcolor=blue,linkcolor=blue]{hyperref}
\setlength{\parindent}{0pt}
\setlength{\parskip}{6pt}
\title{The \texttt{chalkboard} Beamer color theme}
\author{Sewoong Lee\\\small\href{mailto:samuel27@illinois.edu}{samuel27@illinois.edu}}
\date{Version 0.1.0\quad 2026-10-15}
\begin{document}
\maketitle
\section{Overview}
Chalkboard is a Beamer color theme with a dark green background, warm
 off-white text, and yellow, blue, and pink accents. It is distributed as
\texttt{beamercolorthemechalkboard.sty}; the user-facing theme name is
\texttt{chalkboard}. It requires Beamer and no extra fonts or image assets.
\section{Installation and use}
For a single presentation, copy the style file into the directory containing
your main source. On Overleaf, upload it beside the main source. For a local
installation shared between projects, use the directory
\texttt{tex/latex/beamercolorthemechalkboard/} under your personal TeX tree;
find that tree with \texttt{kpsewhich -var-value=TEXMFHOME}.

Load the main theme first and the color theme afterwards:
\begin{verbatim}
\documentclass{beamer}
\usetheme{Madrid}
\usecolortheme{chalkboard}
\begin{document}
\begin{frame}{Hello}
  Normal text and \alert{an important detail}.
\end{frame}
\end{document}
\end{verbatim}
No package options are defined. See \texttt{chalkboard-example.tex} and its
PDF for a title page, lists, mathematics, three block types, and the palette.
Compile either documentation or example source twice with pdfLaTeX from
the package directory.
\newpage
\section{Palette and behavior}
\begin{center}
\begin{tabular}{lll}
\toprule
Color name & HTML value & Main use\\
\midrule
\texttt{chalkboard} & \texttt{2A3B33} & Background\\
\texttt{chalkboarddark} & \texttt{1F2B25} & Titles and palettes\\
\texttt{chalkboardlight} & \texttt{34483E} & Block bodies\\
\texttt{chalkwhite} & \texttt{E6E4D8} & Normal text\\
\texttt{chalkyellow} & \texttt{EBD58A} & Structure and standard titles\\
\texttt{chalkblue} & \texttt{9FC5E0} & Example text and titles\\
\texttt{chalkpink} & \texttt{F0A6A0} & Alerts and alert titles\\
\bottomrule
\end{tabular}
\end{center}
The theme sets normal text and background colors, the four palettes,
title and frame-title colors, author/date colors in headers and footers,
and standard, alert, and example block colors. Projected item labels use
dark green on yellow and a bold font.

In addition to colors, the theme selects rounded blocks without shadows
and the default rounded title page without shadows. It does not select
ball-shaped list markers: these depend on the main or inner theme. To
change these choices, apply your template or font settings after loading
Chalkboard. For example:
\begin{verbatim}
\setbeamertemplate{blocks}[default]
\setbeamerfont{item projected}{series=\mdseries}
\end{verbatim}
\section{Compatibility and customization}
Madrid is the reference combination. The default Beamer theme is also
supported; other themes can supply templates or later settings that affect
the result. The named colors above are available for use in your content.
Use Beamer's \verb|\setbeamercolor| after loading the theme to customize
individual elements.

Handout mode retains the dark background. There is no light/print mode.
Check the result on your presentation equipment: projector brightness and
room lighting can affect readability.
\section{License, maintenance, and history}
Copyright 2026 Sewoong Lee. This work is distributed under the LaTeX Project
Public License version 1.3c or, at your option, any later version.
See \url{https://www.latex-project.org/lppl.txt} and the accompanying
\texttt{LICENSE} for the file list. Maintenance status: maintained.
Current Maintainer: Sewoong Lee, \href{mailto:samuel27@illinois.edu}{samuel27@illinois.edu}.

Version 0.1.0 (2026-10-15): Initial packaged release.

Author profile: \url{https://github.com/SewoongLee}.
\end{document}
